% Proposed replacement for \subsection*{Covariate adjustment} in the Overleaf main.tex (29 October 2026), checked
% against the Snow-polygon primary (primary_snow520_merge_dhsmics_gamma2_v12, with 2001 political stability
% interpolated). main.tex is not edited; paste by hand. The first paragraph is unchanged.
\subsection*{Covariate adjustment}

Covariate selection was guided by the simplified causal processes shown in Figure~\ref{fig:dag}. This causal model was used to guide the selection of variables that are common causes of transmission intensity and mortality, that lie off the causal pathway. These are primarily variables that are proxies for socio-economic status and general health-system quality. We also selected variables that are determinants of childhood mortality from other causes. These include vaccination coverage, child HIV incidence and delivery in a health facility. We did not adjust for LLIN coverage (measured in DHS/MIS surveys) because this is a determinant of prevalence (LLINs reduce transmission) and corresponding reductions in mortality are mediated by reduction in transmission. For malnutrition, causality could go both ways and so measures of malnutrition could mediate the relationship between \textit{Pf}PR$_{2-10}$ and death, or confound it. To assess the potential impact of this, we carried out a sensitivity analysis with and without measures of malnutrition (wasting and stunting).

Adjustment included 17 covariates: 13 area-level summaries computed for each survey unit (see Data sources) and 4 national measures. The area-level covariates were (means or percentages): maternal age at first birth (years); maternal education (years); household wealth quintile (1--5); percentage urban; DTP3 coverage; measles vaccination coverage; percentage of live births delivered in a health facility; the percentage of non-first births with a preceding birth interval shorter than 24 months; household access to improved drinking water; access to improved sanitation; households with electricity; percentage of children with wasting; percentage with stunting. The four maternal and household summaries are means over mothers with a live birth in the 60 months before interview; wealth quintiles are defined within each survey, so the mean wealth quintile measures relative position within the country. National covariates comprised child HIV incidence at ages 0--14 years; GDP per capita; health expenditure per capita; political stability, matched to the calendar year of entry into each age band. HIV incidence, GDP and health expenditure were natural-log transformed. All covariates were centred and scaled by their mean and standard deviation over the analysis sample, and entered linearly, with separate coefficients in each age-band model. The area-level summaries refer to the survey (household conditions and anthropometry at interview; births, deliveries and vaccinations in the preceding one to five years) and take the same value for every age band of that survey, whereas the national covariates vary with the year of band entry.

DHS recode-derived summaries used the women's sampling weights for Births Recode variables and the household weights for household and anthropometric variables. MICS summaries were computed with the MICS women's, household and children's sampling weights. Summaries were calculated from all respondents in each survey unit, before complete-case selection. Questions not asked in a survey were excluded from the denominator; following the DHS tabulation rules, missing or `don't know' answers on vaccination and place of delivery were counted as not vaccinated and not delivered in a facility, and other missing responses were excluded.

Missing DTP3 and measles coverage was replaced by the corresponding WHO/UNICEF national estimate (WUENIC, 2025 revision) for the survey year (27 of the 1,452 survey units in the analysis). Remaining missing area-level covariates were replaced by the mean of available values within the same survey, weighted by the number of records in each unit (one survey unit, for the five household and anthropometric variables). Child HIV incidence was reported for 35 of the 37 countries. For Nigeria, which has no child series, it was imputed with a Bayesian model of log incidence informed by adolescent (15--19 years) incidence and temporal and regional patterns, using the posterior medians; for Liberia, it was derived from UNAIDS counts of new infections in children. For years before South Sudan's independence in 2011, South Sudan took Sudan's GDP per capita, health expenditure and political stability. The Worldwide Governance Indicators were published every two years until 2002, so political stability in 2001 was the mean of each country's 2000 and 2002 values. Records with remaining missing required covariates were excluded (829,521 records and 13,435 deaths). Imputation uncertainty was not propagated in the primary analysis.
